% Generated draft: requires mathematical and source review.
\documentclass[11pt]{article}
\usepackage[margin=25mm]{geometry}
\usepackage{fontspec}
\setmainfont{Latin Modern Roman}
\usepackage{amsmath,amssymb,mathtools,unicode-math}
\setmathfont{Latin Modern Math}
\usepackage{enumitem,needspace,xurl,hyperref}
\setlength{\parindent}{0pt}
\setlength{\parskip}{5pt}
\setlength{\emergencystretch}{3em}
\newcommand{\N}{\mathbb{N}}
\newcommand{\Z}{\mathbb{Z}}
\newcommand{\Q}{\mathbb{Q}}
\newcommand{\R}{\mathbb{R}}
\newcommand{\C}{\mathbb{C}}
\newcommand{\F}{\mathbb{F}}
\newcommand{\A}{\mathbb{A}}
\newcommand{\PP}{\mathbb{P}}
\newcommand{\E}{\mathbb{E}}
\newcommand{\eps}{\varepsilon}
\newcommand{\dd}{\,\mathrm{d}}
\begin{document}
\section*{UnsolvedMath ID: 30000142}
\textbf{Source problem code:} catalogue entry
\subsection{Definitions and mathematical statement}
Let \(M^{2n-1}\) be a finite-volume hyperbolic manifold in the source's intended class. In the usual convention this means a complete manifold with sectional curvature \(-1\), and \(\operatorname{vol}(M)\) denotes its Riemannian volume. The examples in the supplied assessment concern \(n=2,3\), that is dimensions three and five; the exact lower range for the general assertion is not stated.

Write \(B_n(F)\) for the source's intended weight-\(n\) Bloch-type group over a field \(F\), \(\delta_n:B_n(F)\to T_n(F)\) for its boundary map, and \(P_n\) for its intended polylogarithmic volume or regulator functional. The record gives no definitions of these objects, the target \(T_n(F)\), the field, or the functional's normalization. With those source data fixed, the displayed conjecture is
\[
\forall M^{2n-1}\quad\exists r<\infty\quad\exists z_1,\ldots,z_r\in B_n(F)
\quad\text{such that}\quad
\delta_n\left(\sum_{i=1}^r z_i\right)=0,
\qquad
\operatorname{vol}(M)=\sum_{i=1}^r P_n(z_i).
\]
If \(P_n\) is additive on the relevant group, this says that the volume is the regulator of a cycle \(z\in\ker\delta_n\), represented by the indicated finite sum. The condition is on the sum; it does not assert \(\delta_n(z_i)=0\) term by term. No extra rational coefficients, universal prefactor, or power of \(\pi\) is inserted into the source equality.

Bloch-group terminology varies: sometimes a group of symbols is distinguished from the kernel of its boundary map. The source's use of both ``Bloch group'' and a cycle condition must therefore be clarified, rather than making the condition redundant by a guessed definition. Likewise, identifying \(z_i\) as complex numbers rather than group elements would require a source-specified symbol convention. Without these definitions, the displayed equality is an exact account of the proposed form, but not a fully self-contained regulator theorem.

The supplied assessment checked 2026-08-21 reports the hyperbolic-volume formulation established in dimensions three and five in the report, notes advances in related Zagier cases, and says that the general odd-dimensional assertion remains conjectural. No implication from a related weight-four result to all seven-dimensional manifolds is asserted.
\subsection{Short English statement}
Ask whether every intended odd-dimensional finite-volume hyperbolic volume is a finite sum of polylogarithmic values whose combined Bloch boundary vanishes. The Bloch group, base field, regulator normalization, and manifold hypotheses need recovery from the source.
\subsection{Sources}
\begin{itemize}
\item Local UnsolvedMath registry, permanent ID 30000142, catalogue entry, \emph{Bloch-Group Formulas for Odd-Dimensional Hyperbolic Volumes}.
\item \emph{Polylogarithms} (2005), Oberwolfach Reports, pp. 2539--2590, DOI \url{https://doi.org/10.4171/owr/2004/48}.
\item Supplied reference labeled ``weight-four Zagier conjecture'': \url{https://arxiv.org/abs/1803.08585}; no author list or exact theorem is reproduced.
\item Supplied literature triage checked 2026-08-21 and supplied direct statement-verification claim; these are not independently repeated here.
\end{itemize}
\end{document}
